%% figures/opsem.tex — LaCaDiLE operational semantics (Phase 1 T3)
%% Input from main.tex via %% figures/opsem.tex — LaCaDiLE operational semantics (Phase 1 T3)
%% Input from main.tex via %% figures/opsem.tex — LaCaDiLE operational semantics (Phase 1 T3)
%% Input from main.tex via %% figures/opsem.tex — LaCaDiLE operational semantics (Phase 1 T3)
%% Input from main.tex via \input{figures/opsem.tex}
%%
%% Small-step reduction relation: ⟨σ, e⟩ ↦ ⟨σ', e'⟩
%%   σ  — store (Loc → TensorVal)
%%   e  — term
%% Call-by-value, left-to-right evaluation; congruence via evaluation contexts E[·].

\providecommand{\step}[2]{#1 \longmapsto #2}
\providecommand{\config}[2]{\langle #1, #2 \rangle}
\providecommand{\subst}[3]{#1[#2 / #3]}
\providecommand{\fresh}[1]{#1 \text{ fresh}}

\begin{figure}[t]
\small
\begin{mathpar}
\inferrule*[right=\textsc{E-Ctx}]
  {\step{\config{\sigma}{e}}{\config{\sigma'}{e'}}}
  {\step{\config{\sigma}{E[e]}}{\config{\sigma'}{E[e']}}}

\inferrule*[right=\textsc{E-Beta}]
  { }
  {\step{\config{\sigma}{(\lambda x{:}\tau.\, e)\, v}}{\config{\sigma}{\subst{e}{v}{x}}}}

\inferrule*[right=\textsc{E-Let}]
  { }
  {\step{\config{\sigma}{\letin{x = v}{e}}}{\config{\sigma}{\subst{e}{v}{x}}}}

\inferrule*[right=\textsc{E-LetPair}]
  { }
  {\step{\config{\sigma}{\letpair{x}{x'}{(v_1, v_2)}{e}}}{\config{\sigma}{\subst{\subst{e}{v_1}{x}}{v_2}{x'}}}}

\inferrule*[right=\textsc{E-Fst}]
  { }
  {\step{\config{\sigma}{\mathsf{fst}((v_1, v_2))}}{\config{\sigma}{v_1}}}

\inferrule*[right=\textsc{E-Snd}]
  { }
  {\step{\config{\sigma}{\mathsf{snd}((v_1, v_2))}}{\config{\sigma}{v_2}}}
\end{mathpar}
\caption{Core reduction rules. $\textsc{E-Ctx}$ is the congruence rule that
lifts inner-redex reductions to outer terms via evaluation contexts $E[\cdot]$
from the syntax figure. Substitution $\subst{e}{v}{x}$ is capture-avoiding.}
\label{fig:opsem-core}
\end{figure}

\begin{figure}[t]
\small
\begin{mathpar}
\inferrule*[right=\textsc{E-Const}]
  {\fresh{\ell}}
  {\step{\config{\sigma}{\mathsf{const}(v, \bar{d})}}{\config{\sigma[\ell \mapsto (\bar{d}, v)]}{\ell}}}

\inferrule*[right=\textsc{E-Copy}]
  {\fresh{\ell'}}
  {\step{\config{\sigma}{\mathsf{copy}(\ell)}}{\config{\sigma[\ell' \mapsto \sigma(\ell)]}{(\ell, \ell')}}}

\inferrule*[right=\textsc{E-Add}]
  {w = \sigma(\ell_1) \oplus \sigma(\ell_2) \\ \fresh{\ell}}
  {\step{\config{\sigma}{\mathsf{add}(\ell_1, \ell_2)}}{\config{\sigma \setminus \{\ell_1, \ell_2\}\,[\ell \mapsto w]}{\ell}}}

\inferrule*[right=\textsc{E-Mul}]
  {w = \sigma(\ell_1) \odot \sigma(\ell_2) \\ \fresh{\ell}}
  {\step{\config{\sigma}{\mathsf{mul}(\ell_1, \ell_2)}}{\config{\sigma \setminus \{\ell_1, \ell_2\}\,[\ell \mapsto w]}{\ell}}}

\inferrule*[right=\textsc{E-Sum}]
  {w = \mathsf{sum\_val}(\sigma(\ell), i) \\ \fresh{\ell'}}
  {\step{\config{\sigma}{\mathsf{sum}(\ell, i)}}{\config{\sigma \setminus \{\ell\}\,[\ell' \mapsto w]}{\ell'}}}

\inferrule*[right=\textsc{E-Expand}]
  {w = \mathsf{expand\_val}(\sigma(\ell), i, k) \\ \fresh{\ell'}}
  {\step{\config{\sigma}{\mathsf{expand}(\ell, i, k)}}{\config{\sigma \setminus \{\ell\}\,[\ell' \mapsto w]}{\ell'}}}

\inferrule*[right=\textsc{E-UniformLike}]
  {w = \mathsf{uniform\_val}(\sigma(\ell), lo, hi) \\ \fresh{\ell'}}
  {\step{\config{\sigma}{\mathsf{uniform\_like}(\ell, lo, hi)}}{\config{\sigma \setminus \{\ell\}\,[\ell' \mapsto w]}{\ell'}}}
\end{mathpar}
\caption{RISC primitive reduction rules. Each rule evaluates to a location
value after reducing its operands to location values via $\textsc{E-Ctx}$.
$\textsc{E-Const}$ allocates a fresh location for the literal. $\textsc{E-Copy}$
allocates a fresh location $\ell'$ with a physical copy of $\sigma(\ell)$ and
returns a linear pair $(\ell, \ell')$ — the original $\ell$ is retained and
re-bound by the destructuring $\letpair{\cdot}{\cdot}{\cdot}{\cdot}$ in the
caller, while $\ell'$ owns the new location. $\textsc{E-Add}$ and
$\textsc{E-Mul}$ consume both operand locations (removed from $\sigma$) and
allocate a fresh result location. $\textsc{E-Sum}$ and $\textsc{E-Expand}$
consume the single operand and allocate a fresh result location with the
shape-transformed data. $\textsc{E-UniformLike}$ consumes the template location
$\ell$ (freed from $\sigma$) and allocates a fresh location $\ell'$ holding a
new tensor of random values whose shape is read from the template before it is
freed; to retain the template for later use, callers must insert an explicit
$\letpair{t}{t'}{\mathsf{copy}(t)}{\dots}$ before the $\mathsf{uniform\_like}$
call (see T0 \S2.6). Operations $\oplus$, $\odot$, $\mathsf{sum\_val}$,
$\mathsf{expand\_val}$, and $\mathsf{uniform\_val}$ are the underlying
tensor-data operations, defined pointwise on $\mathsf{TensorVal}$ and not
exercised by the metatheory.}
\label{fig:opsem-prims}
\end{figure}

\begin{figure}[t]
\small
\begin{mathpar}
\inferrule*[right=\textsc{E-Handle-Ret}]
  { }
  {\step{\config{\sigma}{\handlewith{\varepsilon_h}{v}{h}}}{\config{\sigma}{v}}}

\inferrule*[right=\textsc{E-Handle-Op}]
  {\mathsf{op}_i \in \varepsilon_h \\
   h = \{\mathsf{op}_j(x_j, k_j) \to e_j\}_j}
  {\step{\config{\sigma}{\handlewith{\varepsilon_h}{E[\performop{op}_i(v)]}{h}}}
        {\config{\sigma}{\subst{\subst{e_i}{v}{x_i}}{K}{k_i}}}}

\\
\text{where } K = \lambda y.\, \handlewith{\varepsilon_h}{E[y]}{h}
\quad\text{(the captured one-shot continuation).}
\end{mathpar}
\caption{Effect handler reduction. $\textsc{E-Handle-Ret}$ discards the handler
when the body has reduced to a value. $\textsc{E-Handle-Op}$ fires when the
body performs an operation in $\varepsilon_h$: the operation's argument $v$ is
substituted into the matching clause $e_i$ for $x_i$, and the captured
continuation $K$ (the evaluation context $E[\cdot]$ around the
$\mathsf{perform}$, re-wrapped by the same handler so nested operations still
see the outer clauses) is substituted for $k_i$. Because $k_i$ is a linear
binding in the handler clause (per \textsc{T-Handle}), every call to $K$
consumes it — the one-shot restriction is mechanical.}
\label{fig:opsem-effects}
\end{figure}

\begin{figure}[t]
\small
\begin{mathpar}
\inferrule*[right=\textsc{E-Grad}]
  { }
  {\step{\config{\sigma}{\mathsf{grad}(\lambda x{:}\tau.\, e)}}
        {\config{\sigma}{\lambda x{:}\tau.\, \lambda g_s{:}\tau_{\mathsf{out}}.\, \handlewith{\{\mathsf{Accum}\}}{\mathsf{adjoint}(e, x, g_s)}{h_{\mathsf{accum}}}}}}

\inferrule*[right=\textsc{E-Vmap}]
  { }
  {\step{\config{\sigma}{\mathsf{vmap}(\lambda x{:}\tau.\, e)}}
        {\config{\sigma}{\lambda x{:}\addDim{d}{\tau}.\, \mathsf{addDimTerm}(d, e)}}}
\end{mathpar}
\[
h_{\mathsf{accum}} \;\triangleq\; \{\mathsf{accum}(p, k) \to k(\mathsf{update\_origin\_buffer}(p))\}
\quad\text{where } p = (\ell, v)
\]
\caption{AD and vectorization transforms. $\textsc{E-Grad}$ wraps the function
body in a two-argument $\lambda$ (parameter $x$ plus seed gradient $g_s$) and
an explicit $\mathsf{handle}[\{\mathsf{Accum}\}]$ whose body is the adjoint
transformation $\mathsf{adjoint}(e, x, g_s)$ of the original function body. The
handler clause routes every $\performop{accum}$ contribution to
$\mathsf{update\_origin\_buffer}$, which keys the gradient buffer by
compile-time origin (tracking $\mathsf{copy}$ chains back to source bindings
and discarding contributions whose origin is a $\mathsf{const}$ literal). The
final value returned by the handler is the gradient buffer entry for $x$'s
origin. Both $\mathsf{adjoint}(e, x, g_s)$ and $\mathsf{addDimTerm}(d, e)$ are
meta-functions on syntax, formalized in Lean during Phase 1 T8; Phase 1's
paper surface treats them as abstract term-to-term transformations whose
defining equations are tabulated in the T0 design document.}
\label{fig:opsem-transforms}
\end{figure}

\noindent\textbf{Meta-function $\mathsf{adjoint}(e, x, g_s)$.}
This is the AD transformation defined informally in the T0 design document
and formalized as a Lean function in Phase 1 T8.
Intuitively, $\mathsf{adjoint}$ walks $e$ in source order and inserts a tape
copy $\letpair{a}{a'}{\mathsf{copy}(a)}{\dots}$ before every primitive whose
adjoint rule requires operand values (notably $\mathsf{mul}$), then walks $e$
in reverse to produce a backward pass that consumes the tape entries, emits
$\performop{accum}$ calls to contribute gradients back to each operand's
origin, and returns the seed gradient $g_s$'s flow through the whole chain as
the final per-parameter gradient. Restrictions: $\mathsf{adjoint}$ is defined
only for $e$ composed of RISC primitives, $\mathsf{let}$, $\mathsf{copy}$,
$\mathsf{letpair}$, pair constructors, $\mathsf{handle}$ with
$\varepsilon_h \subseteq \DiffCompat$, and $\performop{op}$ for
$\mathsf{op} \in \DiffCompat$. It is explicitly undefined on $\mathsf{grad}$
or $\mathsf{vmap}$ subterms inside $e$ — see T0 \S0 for the scope note.

%%
%% Small-step reduction relation: ⟨σ, e⟩ ↦ ⟨σ', e'⟩
%%   σ  — store (Loc → TensorVal)
%%   e  — term
%% Call-by-value, left-to-right evaluation; congruence via evaluation contexts E[·].

\providecommand{\step}[2]{#1 \longmapsto #2}
\providecommand{\config}[2]{\langle #1, #2 \rangle}
\providecommand{\subst}[3]{#1[#2 / #3]}
\providecommand{\fresh}[1]{#1 \text{ fresh}}

\begin{figure}[t]
\small
\begin{mathpar}
\inferrule*[right=\textsc{E-Ctx}]
  {\step{\config{\sigma}{e}}{\config{\sigma'}{e'}}}
  {\step{\config{\sigma}{E[e]}}{\config{\sigma'}{E[e']}}}

\inferrule*[right=\textsc{E-Beta}]
  { }
  {\step{\config{\sigma}{(\lambda x{:}\tau.\, e)\, v}}{\config{\sigma}{\subst{e}{v}{x}}}}

\inferrule*[right=\textsc{E-Let}]
  { }
  {\step{\config{\sigma}{\letin{x = v}{e}}}{\config{\sigma}{\subst{e}{v}{x}}}}

\inferrule*[right=\textsc{E-LetPair}]
  { }
  {\step{\config{\sigma}{\letpair{x}{x'}{(v_1, v_2)}{e}}}{\config{\sigma}{\subst{\subst{e}{v_1}{x}}{v_2}{x'}}}}

\inferrule*[right=\textsc{E-Fst}]
  { }
  {\step{\config{\sigma}{\mathsf{fst}((v_1, v_2))}}{\config{\sigma}{v_1}}}

\inferrule*[right=\textsc{E-Snd}]
  { }
  {\step{\config{\sigma}{\mathsf{snd}((v_1, v_2))}}{\config{\sigma}{v_2}}}
\end{mathpar}
\caption{Core reduction rules. $\textsc{E-Ctx}$ is the congruence rule that
lifts inner-redex reductions to outer terms via evaluation contexts $E[\cdot]$
from the syntax figure. Substitution $\subst{e}{v}{x}$ is capture-avoiding.}
\label{fig:opsem-core}
\end{figure}

\begin{figure}[t]
\small
\begin{mathpar}
\inferrule*[right=\textsc{E-Const}]
  {\fresh{\ell}}
  {\step{\config{\sigma}{\mathsf{const}(v, \bar{d})}}{\config{\sigma[\ell \mapsto (\bar{d}, v)]}{\ell}}}

\inferrule*[right=\textsc{E-Copy}]
  {\fresh{\ell'}}
  {\step{\config{\sigma}{\mathsf{copy}(\ell)}}{\config{\sigma[\ell' \mapsto \sigma(\ell)]}{(\ell, \ell')}}}

\inferrule*[right=\textsc{E-Add}]
  {w = \sigma(\ell_1) \oplus \sigma(\ell_2) \\ \fresh{\ell}}
  {\step{\config{\sigma}{\mathsf{add}(\ell_1, \ell_2)}}{\config{\sigma \setminus \{\ell_1, \ell_2\}\,[\ell \mapsto w]}{\ell}}}

\inferrule*[right=\textsc{E-Mul}]
  {w = \sigma(\ell_1) \odot \sigma(\ell_2) \\ \fresh{\ell}}
  {\step{\config{\sigma}{\mathsf{mul}(\ell_1, \ell_2)}}{\config{\sigma \setminus \{\ell_1, \ell_2\}\,[\ell \mapsto w]}{\ell}}}

\inferrule*[right=\textsc{E-Sum}]
  {w = \mathsf{sum\_val}(\sigma(\ell), i) \\ \fresh{\ell'}}
  {\step{\config{\sigma}{\mathsf{sum}(\ell, i)}}{\config{\sigma \setminus \{\ell\}\,[\ell' \mapsto w]}{\ell'}}}

\inferrule*[right=\textsc{E-Expand}]
  {w = \mathsf{expand\_val}(\sigma(\ell), i, k) \\ \fresh{\ell'}}
  {\step{\config{\sigma}{\mathsf{expand}(\ell, i, k)}}{\config{\sigma \setminus \{\ell\}\,[\ell' \mapsto w]}{\ell'}}}

\inferrule*[right=\textsc{E-UniformLike}]
  {w = \mathsf{uniform\_val}(\sigma(\ell), lo, hi) \\ \fresh{\ell'}}
  {\step{\config{\sigma}{\mathsf{uniform\_like}(\ell, lo, hi)}}{\config{\sigma \setminus \{\ell\}\,[\ell' \mapsto w]}{\ell'}}}
\end{mathpar}
\caption{RISC primitive reduction rules. Each rule evaluates to a location
value after reducing its operands to location values via $\textsc{E-Ctx}$.
$\textsc{E-Const}$ allocates a fresh location for the literal. $\textsc{E-Copy}$
allocates a fresh location $\ell'$ with a physical copy of $\sigma(\ell)$ and
returns a linear pair $(\ell, \ell')$ — the original $\ell$ is retained and
re-bound by the destructuring $\letpair{\cdot}{\cdot}{\cdot}{\cdot}$ in the
caller, while $\ell'$ owns the new location. $\textsc{E-Add}$ and
$\textsc{E-Mul}$ consume both operand locations (removed from $\sigma$) and
allocate a fresh result location. $\textsc{E-Sum}$ and $\textsc{E-Expand}$
consume the single operand and allocate a fresh result location with the
shape-transformed data. $\textsc{E-UniformLike}$ consumes the template location
$\ell$ (freed from $\sigma$) and allocates a fresh location $\ell'$ holding a
new tensor of random values whose shape is read from the template before it is
freed; to retain the template for later use, callers must insert an explicit
$\letpair{t}{t'}{\mathsf{copy}(t)}{\dots}$ before the $\mathsf{uniform\_like}$
call (see T0 \S2.6). Operations $\oplus$, $\odot$, $\mathsf{sum\_val}$,
$\mathsf{expand\_val}$, and $\mathsf{uniform\_val}$ are the underlying
tensor-data operations, defined pointwise on $\mathsf{TensorVal}$ and not
exercised by the metatheory.}
\label{fig:opsem-prims}
\end{figure}

\begin{figure}[t]
\small
\begin{mathpar}
\inferrule*[right=\textsc{E-Handle-Ret}]
  { }
  {\step{\config{\sigma}{\handlewith{\varepsilon_h}{v}{h}}}{\config{\sigma}{v}}}

\inferrule*[right=\textsc{E-Handle-Op}]
  {\mathsf{op}_i \in \varepsilon_h \\
   h = \{\mathsf{op}_j(x_j, k_j) \to e_j\}_j}
  {\step{\config{\sigma}{\handlewith{\varepsilon_h}{E[\performop{op}_i(v)]}{h}}}
        {\config{\sigma}{\subst{\subst{e_i}{v}{x_i}}{K}{k_i}}}}

\\
\text{where } K = \lambda y.\, \handlewith{\varepsilon_h}{E[y]}{h}
\quad\text{(the captured one-shot continuation).}
\end{mathpar}
\caption{Effect handler reduction. $\textsc{E-Handle-Ret}$ discards the handler
when the body has reduced to a value. $\textsc{E-Handle-Op}$ fires when the
body performs an operation in $\varepsilon_h$: the operation's argument $v$ is
substituted into the matching clause $e_i$ for $x_i$, and the captured
continuation $K$ (the evaluation context $E[\cdot]$ around the
$\mathsf{perform}$, re-wrapped by the same handler so nested operations still
see the outer clauses) is substituted for $k_i$. Because $k_i$ is a linear
binding in the handler clause (per \textsc{T-Handle}), every call to $K$
consumes it — the one-shot restriction is mechanical.}
\label{fig:opsem-effects}
\end{figure}

\begin{figure}[t]
\small
\begin{mathpar}
\inferrule*[right=\textsc{E-Grad}]
  { }
  {\step{\config{\sigma}{\mathsf{grad}(\lambda x{:}\tau.\, e)}}
        {\config{\sigma}{\lambda x{:}\tau.\, \lambda g_s{:}\tau_{\mathsf{out}}.\, \handlewith{\{\mathsf{Accum}\}}{\mathsf{adjoint}(e, x, g_s)}{h_{\mathsf{accum}}}}}}

\inferrule*[right=\textsc{E-Vmap}]
  { }
  {\step{\config{\sigma}{\mathsf{vmap}(\lambda x{:}\tau.\, e)}}
        {\config{\sigma}{\lambda x{:}\addDim{d}{\tau}.\, \mathsf{addDimTerm}(d, e)}}}
\end{mathpar}
\[
h_{\mathsf{accum}} \;\triangleq\; \{\mathsf{accum}(p, k) \to k(\mathsf{update\_origin\_buffer}(p))\}
\quad\text{where } p = (\ell, v)
\]
\caption{AD and vectorization transforms. $\textsc{E-Grad}$ wraps the function
body in a two-argument $\lambda$ (parameter $x$ plus seed gradient $g_s$) and
an explicit $\mathsf{handle}[\{\mathsf{Accum}\}]$ whose body is the adjoint
transformation $\mathsf{adjoint}(e, x, g_s)$ of the original function body. The
handler clause routes every $\performop{accum}$ contribution to
$\mathsf{update\_origin\_buffer}$, which keys the gradient buffer by
compile-time origin (tracking $\mathsf{copy}$ chains back to source bindings
and discarding contributions whose origin is a $\mathsf{const}$ literal). The
final value returned by the handler is the gradient buffer entry for $x$'s
origin. Both $\mathsf{adjoint}(e, x, g_s)$ and $\mathsf{addDimTerm}(d, e)$ are
meta-functions on syntax, formalized in Lean during Phase 1 T8; Phase 1's
paper surface treats them as abstract term-to-term transformations whose
defining equations are tabulated in the T0 design document.}
\label{fig:opsem-transforms}
\end{figure}

\noindent\textbf{Meta-function $\mathsf{adjoint}(e, x, g_s)$.}
This is the AD transformation defined informally in the T0 design document
and formalized as a Lean function in Phase 1 T8.
Intuitively, $\mathsf{adjoint}$ walks $e$ in source order and inserts a tape
copy $\letpair{a}{a'}{\mathsf{copy}(a)}{\dots}$ before every primitive whose
adjoint rule requires operand values (notably $\mathsf{mul}$), then walks $e$
in reverse to produce a backward pass that consumes the tape entries, emits
$\performop{accum}$ calls to contribute gradients back to each operand's
origin, and returns the seed gradient $g_s$'s flow through the whole chain as
the final per-parameter gradient. Restrictions: $\mathsf{adjoint}$ is defined
only for $e$ composed of RISC primitives, $\mathsf{let}$, $\mathsf{copy}$,
$\mathsf{letpair}$, pair constructors, $\mathsf{handle}$ with
$\varepsilon_h \subseteq \DiffCompat$, and $\performop{op}$ for
$\mathsf{op} \in \DiffCompat$. It is explicitly undefined on $\mathsf{grad}$
or $\mathsf{vmap}$ subterms inside $e$ — see T0 \S0 for the scope note.

%%
%% Small-step reduction relation: ⟨σ, e⟩ ↦ ⟨σ', e'⟩
%%   σ  — store (Loc → TensorVal)
%%   e  — term
%% Call-by-value, left-to-right evaluation; congruence via evaluation contexts E[·].

\providecommand{\step}[2]{#1 \longmapsto #2}
\providecommand{\config}[2]{\langle #1, #2 \rangle}
\providecommand{\subst}[3]{#1[#2 / #3]}
\providecommand{\fresh}[1]{#1 \text{ fresh}}

\begin{figure}[t]
\small
\begin{mathpar}
\inferrule*[right=\textsc{E-Ctx}]
  {\step{\config{\sigma}{e}}{\config{\sigma'}{e'}}}
  {\step{\config{\sigma}{E[e]}}{\config{\sigma'}{E[e']}}}

\inferrule*[right=\textsc{E-Beta}]
  { }
  {\step{\config{\sigma}{(\lambda x{:}\tau.\, e)\, v}}{\config{\sigma}{\subst{e}{v}{x}}}}

\inferrule*[right=\textsc{E-Let}]
  { }
  {\step{\config{\sigma}{\letin{x = v}{e}}}{\config{\sigma}{\subst{e}{v}{x}}}}

\inferrule*[right=\textsc{E-LetPair}]
  { }
  {\step{\config{\sigma}{\letpair{x}{x'}{(v_1, v_2)}{e}}}{\config{\sigma}{\subst{\subst{e}{v_1}{x}}{v_2}{x'}}}}

\inferrule*[right=\textsc{E-Fst}]
  { }
  {\step{\config{\sigma}{\mathsf{fst}((v_1, v_2))}}{\config{\sigma}{v_1}}}

\inferrule*[right=\textsc{E-Snd}]
  { }
  {\step{\config{\sigma}{\mathsf{snd}((v_1, v_2))}}{\config{\sigma}{v_2}}}
\end{mathpar}
\caption{Core reduction rules. $\textsc{E-Ctx}$ is the congruence rule that
lifts inner-redex reductions to outer terms via evaluation contexts $E[\cdot]$
from the syntax figure. Substitution $\subst{e}{v}{x}$ is capture-avoiding.}
\label{fig:opsem-core}
\end{figure}

\begin{figure}[t]
\small
\begin{mathpar}
\inferrule*[right=\textsc{E-Const}]
  {\fresh{\ell}}
  {\step{\config{\sigma}{\mathsf{const}(v, \bar{d})}}{\config{\sigma[\ell \mapsto (\bar{d}, v)]}{\ell}}}

\inferrule*[right=\textsc{E-Copy}]
  {\fresh{\ell'}}
  {\step{\config{\sigma}{\mathsf{copy}(\ell)}}{\config{\sigma[\ell' \mapsto \sigma(\ell)]}{(\ell, \ell')}}}

\inferrule*[right=\textsc{E-Add}]
  {w = \sigma(\ell_1) \oplus \sigma(\ell_2) \\ \fresh{\ell}}
  {\step{\config{\sigma}{\mathsf{add}(\ell_1, \ell_2)}}{\config{\sigma \setminus \{\ell_1, \ell_2\}\,[\ell \mapsto w]}{\ell}}}

\inferrule*[right=\textsc{E-Mul}]
  {w = \sigma(\ell_1) \odot \sigma(\ell_2) \\ \fresh{\ell}}
  {\step{\config{\sigma}{\mathsf{mul}(\ell_1, \ell_2)}}{\config{\sigma \setminus \{\ell_1, \ell_2\}\,[\ell \mapsto w]}{\ell}}}

\inferrule*[right=\textsc{E-Sum}]
  {w = \mathsf{sum\_val}(\sigma(\ell), i) \\ \fresh{\ell'}}
  {\step{\config{\sigma}{\mathsf{sum}(\ell, i)}}{\config{\sigma \setminus \{\ell\}\,[\ell' \mapsto w]}{\ell'}}}

\inferrule*[right=\textsc{E-Expand}]
  {w = \mathsf{expand\_val}(\sigma(\ell), i, k) \\ \fresh{\ell'}}
  {\step{\config{\sigma}{\mathsf{expand}(\ell, i, k)}}{\config{\sigma \setminus \{\ell\}\,[\ell' \mapsto w]}{\ell'}}}

\inferrule*[right=\textsc{E-UniformLike}]
  {w = \mathsf{uniform\_val}(\sigma(\ell), lo, hi) \\ \fresh{\ell'}}
  {\step{\config{\sigma}{\mathsf{uniform\_like}(\ell, lo, hi)}}{\config{\sigma \setminus \{\ell\}\,[\ell' \mapsto w]}{\ell'}}}
\end{mathpar}
\caption{RISC primitive reduction rules. Each rule evaluates to a location
value after reducing its operands to location values via $\textsc{E-Ctx}$.
$\textsc{E-Const}$ allocates a fresh location for the literal. $\textsc{E-Copy}$
allocates a fresh location $\ell'$ with a physical copy of $\sigma(\ell)$ and
returns a linear pair $(\ell, \ell')$ — the original $\ell$ is retained and
re-bound by the destructuring $\letpair{\cdot}{\cdot}{\cdot}{\cdot}$ in the
caller, while $\ell'$ owns the new location. $\textsc{E-Add}$ and
$\textsc{E-Mul}$ consume both operand locations (removed from $\sigma$) and
allocate a fresh result location. $\textsc{E-Sum}$ and $\textsc{E-Expand}$
consume the single operand and allocate a fresh result location with the
shape-transformed data. $\textsc{E-UniformLike}$ consumes the template location
$\ell$ (freed from $\sigma$) and allocates a fresh location $\ell'$ holding a
new tensor of random values whose shape is read from the template before it is
freed; to retain the template for later use, callers must insert an explicit
$\letpair{t}{t'}{\mathsf{copy}(t)}{\dots}$ before the $\mathsf{uniform\_like}$
call (see T0 \S2.6). Operations $\oplus$, $\odot$, $\mathsf{sum\_val}$,
$\mathsf{expand\_val}$, and $\mathsf{uniform\_val}$ are the underlying
tensor-data operations, defined pointwise on $\mathsf{TensorVal}$ and not
exercised by the metatheory.}
\label{fig:opsem-prims}
\end{figure}

\begin{figure}[t]
\small
\begin{mathpar}
\inferrule*[right=\textsc{E-Handle-Ret}]
  { }
  {\step{\config{\sigma}{\handlewith{\varepsilon_h}{v}{h}}}{\config{\sigma}{v}}}

\inferrule*[right=\textsc{E-Handle-Op}]
  {\mathsf{op}_i \in \varepsilon_h \\
   h = \{\mathsf{op}_j(x_j, k_j) \to e_j\}_j}
  {\step{\config{\sigma}{\handlewith{\varepsilon_h}{E[\performop{op}_i(v)]}{h}}}
        {\config{\sigma}{\subst{\subst{e_i}{v}{x_i}}{K}{k_i}}}}

\\
\text{where } K = \lambda y.\, \handlewith{\varepsilon_h}{E[y]}{h}
\quad\text{(the captured one-shot continuation).}
\end{mathpar}
\caption{Effect handler reduction. $\textsc{E-Handle-Ret}$ discards the handler
when the body has reduced to a value. $\textsc{E-Handle-Op}$ fires when the
body performs an operation in $\varepsilon_h$: the operation's argument $v$ is
substituted into the matching clause $e_i$ for $x_i$, and the captured
continuation $K$ (the evaluation context $E[\cdot]$ around the
$\mathsf{perform}$, re-wrapped by the same handler so nested operations still
see the outer clauses) is substituted for $k_i$. Because $k_i$ is a linear
binding in the handler clause (per \textsc{T-Handle}), every call to $K$
consumes it — the one-shot restriction is mechanical.}
\label{fig:opsem-effects}
\end{figure}

\begin{figure}[t]
\small
\begin{mathpar}
\inferrule*[right=\textsc{E-Grad}]
  { }
  {\step{\config{\sigma}{\mathsf{grad}(\lambda x{:}\tau.\, e)}}
        {\config{\sigma}{\lambda x{:}\tau.\, \lambda g_s{:}\tau_{\mathsf{out}}.\, \handlewith{\{\mathsf{Accum}\}}{\mathsf{adjoint}(e, x, g_s)}{h_{\mathsf{accum}}}}}}

\inferrule*[right=\textsc{E-Vmap}]
  { }
  {\step{\config{\sigma}{\mathsf{vmap}(\lambda x{:}\tau.\, e)}}
        {\config{\sigma}{\lambda x{:}\addDim{d}{\tau}.\, \mathsf{addDimTerm}(d, e)}}}
\end{mathpar}
\[
h_{\mathsf{accum}} \;\triangleq\; \{\mathsf{accum}(p, k) \to k(\mathsf{update\_origin\_buffer}(p))\}
\quad\text{where } p = (\ell, v)
\]
\caption{AD and vectorization transforms. $\textsc{E-Grad}$ wraps the function
body in a two-argument $\lambda$ (parameter $x$ plus seed gradient $g_s$) and
an explicit $\mathsf{handle}[\{\mathsf{Accum}\}]$ whose body is the adjoint
transformation $\mathsf{adjoint}(e, x, g_s)$ of the original function body. The
handler clause routes every $\performop{accum}$ contribution to
$\mathsf{update\_origin\_buffer}$, which keys the gradient buffer by
compile-time origin (tracking $\mathsf{copy}$ chains back to source bindings
and discarding contributions whose origin is a $\mathsf{const}$ literal). The
final value returned by the handler is the gradient buffer entry for $x$'s
origin. Both $\mathsf{adjoint}(e, x, g_s)$ and $\mathsf{addDimTerm}(d, e)$ are
meta-functions on syntax, formalized in Lean during Phase 1 T8; Phase 1's
paper surface treats them as abstract term-to-term transformations whose
defining equations are tabulated in the T0 design document.}
\label{fig:opsem-transforms}
\end{figure}

\noindent\textbf{Meta-function $\mathsf{adjoint}(e, x, g_s)$.}
This is the AD transformation defined informally in the T0 design document
and formalized as a Lean function in Phase 1 T8.
Intuitively, $\mathsf{adjoint}$ walks $e$ in source order and inserts a tape
copy $\letpair{a}{a'}{\mathsf{copy}(a)}{\dots}$ before every primitive whose
adjoint rule requires operand values (notably $\mathsf{mul}$), then walks $e$
in reverse to produce a backward pass that consumes the tape entries, emits
$\performop{accum}$ calls to contribute gradients back to each operand's
origin, and returns the seed gradient $g_s$'s flow through the whole chain as
the final per-parameter gradient. Restrictions: $\mathsf{adjoint}$ is defined
only for $e$ composed of RISC primitives, $\mathsf{let}$, $\mathsf{copy}$,
$\mathsf{letpair}$, pair constructors, $\mathsf{handle}$ with
$\varepsilon_h \subseteq \DiffCompat$, and $\performop{op}$ for
$\mathsf{op} \in \DiffCompat$. It is explicitly undefined on $\mathsf{grad}$
or $\mathsf{vmap}$ subterms inside $e$ — see T0 \S0 for the scope note.

%%
%% Small-step reduction relation: ⟨σ, e⟩ ↦ ⟨σ', e'⟩
%%   σ  — store (Loc → TensorVal)
%%   e  — term
%% Call-by-value, left-to-right evaluation; congruence via evaluation contexts E[·].

\providecommand{\step}[2]{#1 \longmapsto #2}
\providecommand{\config}[2]{\langle #1, #2 \rangle}
\providecommand{\subst}[3]{#1[#2 / #3]}
\providecommand{\fresh}[1]{#1 \text{ fresh}}

\begin{figure}[t]
\small
\begin{mathpar}
\inferrule*[right=\textsc{E-Ctx}]
  {\step{\config{\sigma}{e}}{\config{\sigma'}{e'}}}
  {\step{\config{\sigma}{E[e]}}{\config{\sigma'}{E[e']}}}

\inferrule*[right=\textsc{E-Beta}]
  { }
  {\step{\config{\sigma}{(\lambda x{:}\tau.\, e)\, v}}{\config{\sigma}{\subst{e}{v}{x}}}}

\inferrule*[right=\textsc{E-Let}]
  { }
  {\step{\config{\sigma}{\letin{x = v}{e}}}{\config{\sigma}{\subst{e}{v}{x}}}}

\inferrule*[right=\textsc{E-LetPair}]
  { }
  {\step{\config{\sigma}{\letpair{x}{x'}{(v_1, v_2)}{e}}}{\config{\sigma}{\subst{\subst{e}{v_1}{x}}{v_2}{x'}}}}

\inferrule*[right=\textsc{E-Fst}]
  { }
  {\step{\config{\sigma}{\mathsf{fst}((v_1, v_2))}}{\config{\sigma}{v_1}}}

\inferrule*[right=\textsc{E-Snd}]
  { }
  {\step{\config{\sigma}{\mathsf{snd}((v_1, v_2))}}{\config{\sigma}{v_2}}}
\end{mathpar}
\caption{Core reduction rules. $\textsc{E-Ctx}$ is the congruence rule that
lifts inner-redex reductions to outer terms via evaluation contexts $E[\cdot]$
from the syntax figure. Substitution $\subst{e}{v}{x}$ is capture-avoiding.}
\label{fig:opsem-core}
\end{figure}

\begin{figure}[t]
\small
\begin{mathpar}
\inferrule*[right=\textsc{E-Const}]
  {\fresh{\ell}}
  {\step{\config{\sigma}{\mathsf{const}(v, \bar{d})}}{\config{\sigma[\ell \mapsto (\bar{d}, v)]}{\ell}}}

\inferrule*[right=\textsc{E-Copy}]
  {\fresh{\ell'}}
  {\step{\config{\sigma}{\mathsf{copy}(\ell)}}{\config{\sigma[\ell' \mapsto \sigma(\ell)]}{(\ell, \ell')}}}

\inferrule*[right=\textsc{E-Add}]
  {w = \sigma(\ell_1) \oplus \sigma(\ell_2) \\ \fresh{\ell}}
  {\step{\config{\sigma}{\mathsf{add}(\ell_1, \ell_2)}}{\config{\sigma \setminus \{\ell_1, \ell_2\}\,[\ell \mapsto w]}{\ell}}}

\inferrule*[right=\textsc{E-Mul}]
  {w = \sigma(\ell_1) \odot \sigma(\ell_2) \\ \fresh{\ell}}
  {\step{\config{\sigma}{\mathsf{mul}(\ell_1, \ell_2)}}{\config{\sigma \setminus \{\ell_1, \ell_2\}\,[\ell \mapsto w]}{\ell}}}

\inferrule*[right=\textsc{E-Sum}]
  {w = \mathsf{sum\_val}(\sigma(\ell), i) \\ \fresh{\ell'}}
  {\step{\config{\sigma}{\mathsf{sum}(\ell, i)}}{\config{\sigma \setminus \{\ell\}\,[\ell' \mapsto w]}{\ell'}}}

\inferrule*[right=\textsc{E-Expand}]
  {w = \mathsf{expand\_val}(\sigma(\ell), i, k) \\ \fresh{\ell'}}
  {\step{\config{\sigma}{\mathsf{expand}(\ell, i, k)}}{\config{\sigma \setminus \{\ell\}\,[\ell' \mapsto w]}{\ell'}}}

\inferrule*[right=\textsc{E-UniformLike}]
  {w = \mathsf{uniform\_val}(\sigma(\ell), lo, hi) \\ \fresh{\ell'}}
  {\step{\config{\sigma}{\mathsf{uniform\_like}(\ell, lo, hi)}}{\config{\sigma \setminus \{\ell\}\,[\ell' \mapsto w]}{\ell'}}}
\end{mathpar}
\caption{RISC primitive reduction rules. Each rule evaluates to a location
value after reducing its operands to location values via $\textsc{E-Ctx}$.
$\textsc{E-Const}$ allocates a fresh location for the literal. $\textsc{E-Copy}$
allocates a fresh location $\ell'$ with a physical copy of $\sigma(\ell)$ and
returns a linear pair $(\ell, \ell')$ — the original $\ell$ is retained and
re-bound by the destructuring $\letpair{\cdot}{\cdot}{\cdot}{\cdot}$ in the
caller, while $\ell'$ owns the new location. $\textsc{E-Add}$ and
$\textsc{E-Mul}$ consume both operand locations (removed from $\sigma$) and
allocate a fresh result location. $\textsc{E-Sum}$ and $\textsc{E-Expand}$
consume the single operand and allocate a fresh result location with the
shape-transformed data. $\textsc{E-UniformLike}$ consumes the template location
$\ell$ (freed from $\sigma$) and allocates a fresh location $\ell'$ holding a
new tensor of random values whose shape is read from the template before it is
freed; to retain the template for later use, callers must insert an explicit
$\letpair{t}{t'}{\mathsf{copy}(t)}{\dots}$ before the $\mathsf{uniform\_like}$
call (see T0 \S2.6). Operations $\oplus$, $\odot$, $\mathsf{sum\_val}$,
$\mathsf{expand\_val}$, and $\mathsf{uniform\_val}$ are the underlying
tensor-data operations, defined pointwise on $\mathsf{TensorVal}$ and not
exercised by the metatheory.}
\label{fig:opsem-prims}
\end{figure}

\begin{figure}[t]
\small
\begin{mathpar}
\inferrule*[right=\textsc{E-Handle-Ret}]
  { }
  {\step{\config{\sigma}{\handlewith{\varepsilon_h}{v}{h}}}{\config{\sigma}{v}}}

\inferrule*[right=\textsc{E-Handle-Op}]
  {\mathsf{op}_i \in \varepsilon_h \\
   h = \{\mathsf{op}_j(x_j, k_j) \to e_j\}_j}
  {\step{\config{\sigma}{\handlewith{\varepsilon_h}{E[\performop{op}_i(v)]}{h}}}
        {\config{\sigma}{\subst{\subst{e_i}{v}{x_i}}{K}{k_i}}}}

\\
\text{where } K = \lambda y.\, \handlewith{\varepsilon_h}{E[y]}{h}
\quad\text{(the captured one-shot continuation).}
\end{mathpar}
\caption{Effect handler reduction. $\textsc{E-Handle-Ret}$ discards the handler
when the body has reduced to a value. $\textsc{E-Handle-Op}$ fires when the
body performs an operation in $\varepsilon_h$: the operation's argument $v$ is
substituted into the matching clause $e_i$ for $x_i$, and the captured
continuation $K$ (the evaluation context $E[\cdot]$ around the
$\mathsf{perform}$, re-wrapped by the same handler so nested operations still
see the outer clauses) is substituted for $k_i$. Because $k_i$ is a linear
binding in the handler clause (per \textsc{T-Handle}), every call to $K$
consumes it — the one-shot restriction is mechanical.}
\label{fig:opsem-effects}
\end{figure}

\begin{figure}[t]
\small
\begin{mathpar}
\inferrule*[right=\textsc{E-Grad}]
  { }
  {\step{\config{\sigma}{\mathsf{grad}(\lambda x{:}\tau.\, e)}}
        {\config{\sigma}{\lambda x{:}\tau.\, \lambda g_s{:}\tau_{\mathsf{out}}.\, \handlewith{\{\mathsf{Accum}\}}{\mathsf{adjoint}(e, x, g_s)}{h_{\mathsf{accum}}}}}}

\inferrule*[right=\textsc{E-Vmap}]
  { }
  {\step{\config{\sigma}{\mathsf{vmap}(\lambda x{:}\tau.\, e)}}
        {\config{\sigma}{\lambda x{:}\addDim{d}{\tau}.\, \mathsf{addDimTerm}(d, e)}}}
\end{mathpar}
\[
h_{\mathsf{accum}} \;\triangleq\; \{\mathsf{accum}(p, k) \to k(\mathsf{update\_origin\_buffer}(p))\}
\quad\text{where } p = (\ell, v)
\]
\caption{AD and vectorization transforms. $\textsc{E-Grad}$ wraps the function
body in a two-argument $\lambda$ (parameter $x$ plus seed gradient $g_s$) and
an explicit $\mathsf{handle}[\{\mathsf{Accum}\}]$ whose body is the adjoint
transformation $\mathsf{adjoint}(e, x, g_s)$ of the original function body. The
handler clause routes every $\performop{accum}$ contribution to
$\mathsf{update\_origin\_buffer}$, which keys the gradient buffer by
compile-time origin (tracking $\mathsf{copy}$ chains back to source bindings
and discarding contributions whose origin is a $\mathsf{const}$ literal). The
final value returned by the handler is the gradient buffer entry for $x$'s
origin. Both $\mathsf{adjoint}(e, x, g_s)$ and $\mathsf{addDimTerm}(d, e)$ are
meta-functions on syntax, formalized in Lean during Phase 1 T8; Phase 1's
paper surface treats them as abstract term-to-term transformations whose
defining equations are tabulated in the T0 design document.}
\label{fig:opsem-transforms}
\end{figure}

\noindent\textbf{Meta-function $\mathsf{adjoint}(e, x, g_s)$.}
This is the AD transformation defined informally in the T0 design document
and formalized as a Lean function in Phase 1 T8.
Intuitively, $\mathsf{adjoint}$ walks $e$ in source order and inserts a tape
copy $\letpair{a}{a'}{\mathsf{copy}(a)}{\dots}$ before every primitive whose
adjoint rule requires operand values (notably $\mathsf{mul}$), then walks $e$
in reverse to produce a backward pass that consumes the tape entries, emits
$\performop{accum}$ calls to contribute gradients back to each operand's
origin, and returns the seed gradient $g_s$'s flow through the whole chain as
the final per-parameter gradient. Restrictions: $\mathsf{adjoint}$ is defined
only for $e$ composed of RISC primitives, $\mathsf{let}$, $\mathsf{copy}$,
$\mathsf{letpair}$, pair constructors, $\mathsf{handle}$ with
$\varepsilon_h \subseteq \DiffCompat$, and $\performop{op}$ for
$\mathsf{op} \in \DiffCompat$. It is explicitly undefined on $\mathsf{grad}$
or $\mathsf{vmap}$ subterms inside $e$ — see T0 \S0 for the scope note.
